\chapter{设备管理}

\section{I/O 系统与设备分类}
\concept{设备管理的对象与任务}
\begin{points}
  \item 除 CPU 和内存外，大部分硬件设备称为外部设备。
  \item I/O 系统包括 I/O 设备、接口线路、控制部件、通道和管理软件。
  \item I/O 操作是主存与外围设备介质之间的信息传送。
  \item 设备管理任务：响应 I/O 请求、分配 I/O 设备、管理设备可靠性与安全性、提供友好访问接口。
\end{points}
\plain{I/O 部分就是解决 CPU 快、设备慢、设备复杂且共享的问题。DMA/通道减少 CPU 搬运，缓冲平滑速度差，SPOOLing 把独占设备虚拟成共享设备。}
\exam{常考设备分类、四种 I/O 控制方式、DMA 与中断/通道区别、缓冲和 SPOOLing 的作用。}


\concept{设备管理四个功能}
\begin{points}
  \item 设备分配：根据用户 I/O 请求分配设备、控制器和通道。
  \item 设备处理：启动设备，并在 I/O 完成时处理中断。
  \item 缓冲管理：缓和 CPU 与 I/O 速度不匹配，管理缓冲区分配与释放。
  \item 设备独立性：用户编程使用逻辑设备名，不依赖具体物理设备。
\end{points}
\plain{I/O 部分就是解决 CPU 快、设备慢、设备复杂且共享的问题。DMA/通道减少 CPU 搬运，缓冲平滑速度差，SPOOLing 把独占设备虚拟成共享设备。}
\exam{常考设备分类、四种 I/O 控制方式、DMA 与中断/通道区别、缓冲和 SPOOLing 的作用。}


\concept{按使用特性分类}
\begin{points}
  \item 存储设备：保存信息，如磁盘、磁带、SSD。
  \item 传输设备：传输数据，如网卡、调制解调器。
  \item 人机交互设备：输入或输出用户信息，如键盘、显示器、打印机。
\end{points}
\plain{把这个知识点放回本章主线理解：它回答的是 OS 为了管理资源、隐藏硬件、保证并发正确性而采用的某个对象、规则或步骤。}
\exam{常考选择/填空/简答，让你判断概念归属、比较相近概念，或按“定义-作用-机制-易错点”展开。}


\concept{按共享属性分类}
\begin{points}
  \item 独占设备：一段时间内只允许一个进程使用，如传统打印机。
  \item 共享设备：允许多个进程交替或并发使用，如磁盘。
  \item 虚拟设备：通过虚拟技术把独占设备改造成多个逻辑设备，如 SPOOLing 打印机。
\end{points}
\plain{多个进程都想用同一批资源，所以必须规定能不能一起用、谁先用、用完怎么释放。}
\exam{常考互斥共享/同时共享举例，或者问并发和共享为什么是 OS 最基本特征。}


\concept{按信息交换单位分类}
\begin{points}
  \item 块设备以固定大小数据块为单位传输，如磁盘、磁带。
  \item 字符设备以字符流为单位传输，如键盘、鼠标、串口、打印机。
  \item 选择题中，磁盘通常是块设备；键盘、打印机、显示器通常是字符设备。
\end{points}
\plain{把这个知识点放回本章主线理解：它回答的是 OS 为了管理资源、隐藏硬件、保证并发正确性而采用的某个对象、规则或步骤。}
\exam{常考选择/填空/简答，让你判断概念归属、比较相近概念，或按“定义-作用-机制-易错点”展开。}


\section{I/O 硬件}
\concept{设备控制器与寄存器}
\begin{points}
  \item 设备控制器是 CPU 与外设之间的电子部件，负责控制设备操作。
  \item 控制器通常包含数据寄存器、状态寄存器、控制寄存器。
  \item CPU 通过读写控制器寄存器与设备交互。
  \item 设备完成操作或发生错误时，通过中断通知 CPU。
\end{points}
\plain{硬件基础不用背太深，重点是知道哪些硬件机制支撑 OS：寄存器保存现场，PSW 保存状态，CPU 模式限制权限。}
\exam{常考 CPU 工作流程、单/多处理器、寄存器类别、PSW 中断位/状态位。}


\concept{I/O 地址空间}
\begin{points}
  \item 端口映射 I/O：I/O 设备有独立地址空间，CPU 使用专门 I/O 指令访问端口。
  \item 内存映射 I/O：设备寄存器映射到内存地址空间，CPU 用普通访存指令访问。
  \item 二者都需要保护，普通用户程序不能随意访问设备寄存器。
\end{points}
\plain{程序看到的是逻辑地址，内存条上真正用的是物理地址，中间靠装载器/MMU 转换。}
\exam{常考逻辑地址到物理地址转换、静态/动态重定位、越界检查。}


\concept{通道}
\begin{points}
  \item 通道是一种专门负责 I/O 控制的处理部件，比普通设备控制器更强。
  \item 通道能执行通道程序，组织多个 I/O 操作，减轻 CPU 负担。
  \item 处理机和通道可以并行工作，因此通道提高了大型系统 I/O 并行能力。
\end{points}
\plain{I/O 部分就是解决 CPU 快、设备慢、设备复杂且共享的问题。DMA/通道减少 CPU 搬运，缓冲平滑速度差，SPOOLing 把独占设备虚拟成共享设备。}
\exam{常考设备分类、四种 I/O 控制方式、DMA 与中断/通道区别、缓冲和 SPOOLing 的作用。}


\section{I/O 控制方式}
\concept{程序直接控制方式}
\begin{points}
  \item CPU 反复查询设备状态寄存器，直到设备就绪，再进行数据传送。
  \item 优点：实现简单。
  \item 缺点：CPU 忙等，浪费处理机时间。
  \item 又称轮询方式，适合简单或低速设备。
\end{points}
\plain{把这个知识点放回本章主线理解：它回答的是 OS 为了管理资源、隐藏硬件、保证并发正确性而采用的某个对象、规则或步骤。}
\exam{常考选择/填空/简答，让你判断概念归属、比较相近概念，或按“定义-作用-机制-易错点”展开。}


\concept{中断驱动方式}
\begin{points}
  \item CPU 启动 I/O 后转去执行其他任务，设备完成后发中断通知 CPU。
  \item 优点：减少 CPU 忙等，提高并发性。
  \item 缺点：每次传送仍需 CPU 参与，频繁中断会增加开销。
\end{points}
\plain{中断就是硬件或软件打断 CPU 当前执行，让 OS 先处理更紧急或必须处理的事件。}
\exam{常考中断处理步骤、中断与异常/系统调用区别、为什么中断能提高 CPU 与 I/O 并行度。}


\concept{DMA 方式}
\begin{points}
  \item DMA 允许设备控制器在主存和设备之间直接传送一批数据，无需 CPU 逐字节搬运。
  \item CPU 只负责初始化 DMA 控制器，设置内存地址、传输长度、方向等。
  \item 传输完成后 DMA 控制器发中断通知 CPU。
  \item 优点：大幅减少 CPU 参与，适合磁盘、网卡等高速块设备。
\end{points}
\plain{I/O 部分就是解决 CPU 快、设备慢、设备复杂且共享的问题。DMA/通道减少 CPU 搬运，缓冲平滑速度差，SPOOLing 把独占设备虚拟成共享设备。}
\exam{常考设备分类、四种 I/O 控制方式、DMA 与中断/通道区别、缓冲和 SPOOLing 的作用。}


\concept{通道控制方式}
\begin{points}
  \item CPU 把一组 I/O 操作写成通道程序交给通道执行。
  \item 通道独立控制设备完成复杂 I/O 序列，完成后中断 CPU。
  \item 与 DMA 相比，通道更像专用 I/O 处理机，控制能力更强。
\end{points}
\plain{I/O 部分就是解决 CPU 快、设备慢、设备复杂且共享的问题。DMA/通道减少 CPU 搬运，缓冲平滑速度差，SPOOLing 把独占设备虚拟成共享设备。}
\exam{常考设备分类、四种 I/O 控制方式、DMA 与中断/通道区别、缓冲和 SPOOLing 的作用。}


\section{I/O 软件原理}
\concept{I/O 软件层次}
\begin{points}
  \item 用户层 I/O 软件：库函数、假脱机程序等，为用户提供高级接口。
  \item 设备独立性软件：提供统一逻辑设备接口、命名、保护、缓冲、错误处理。
  \item 设备驱动程序：针对具体设备控制器发命令、处理中断。
  \item 中断处理程序：响应设备中断，完成状态检查、唤醒等待进程等。
\end{points}
\plain{I/O 部分就是解决 CPU 快、设备慢、设备复杂且共享的问题。DMA/通道减少 CPU 搬运，缓冲平滑速度差，SPOOLing 把独占设备虚拟成共享设备。}
\exam{常考设备分类、四种 I/O 控制方式、DMA 与中断/通道区别、缓冲和 SPOOLing 的作用。}


\concept{设备独立性}
\begin{points}
  \item 设备独立性又称设备无关性，指用户程序不依赖具体物理设备。
  \item 第一层含义：独立于同类设备的具体设备号，例如打印到逻辑打印机而非某台具体打印机。
  \item 第二层含义：在一定抽象下独立于设备类型，例如统一用文件接口访问普通文件和部分设备文件。
  \item 好处：提高可移植性、灵活性和设备分配效率。
\end{points}
\plain{I/O 部分就是解决 CPU 快、设备慢、设备复杂且共享的问题。DMA/通道减少 CPU 搬运，缓冲平滑速度差，SPOOLing 把独占设备虚拟成共享设备。}
\exam{常考设备分类、四种 I/O 控制方式、DMA 与中断/通道区别、缓冲和 SPOOLing 的作用。}


\section{缓冲技术管理}
\concept{为什么需要缓冲}
\begin{points}
  \item CPU 和 I/O 设备速度差异巨大，缓冲可缓和二者速度不匹配。
  \item 缓冲可减少中断次数和设备启动次数，提高并行程度。
  \item 缓冲还能适配不同数据粒度，例如用户按字节读写，磁盘按块传输。
\end{points}
\plain{I/O 部分就是解决 CPU 快、设备慢、设备复杂且共享的问题。DMA/通道减少 CPU 搬运，缓冲平滑速度差，SPOOLing 把独占设备虚拟成共享设备。}
\exam{常考设备分类、四种 I/O 控制方式、DMA 与中断/通道区别、缓冲和 SPOOLing 的作用。}


\concept{单缓冲}
\begin{points}
  \item 系统设置一个缓冲区，设备先把数据送入缓冲区，再由进程取走。
  \item 相比无缓冲，可让设备传输和进程处理部分重叠。
  \item 但一个缓冲区仍可能使设备或进程等待。
\end{points}
\plain{I/O 部分就是解决 CPU 快、设备慢、设备复杂且共享的问题。DMA/通道减少 CPU 搬运，缓冲平滑速度差，SPOOLing 把独占设备虚拟成共享设备。}
\exam{常考设备分类、四种 I/O 控制方式、DMA 与中断/通道区别、缓冲和 SPOOLing 的作用。}


\concept{双缓冲}
\begin{points}
  \item 设置两个缓冲区，一个用于设备填充，另一个供进程处理。
  \item 可以更好地实现输入/输出与处理并行。
  \item 常用于数据流连续、生产和消费速度接近的场景。
\end{points}
\plain{I/O 部分就是解决 CPU 快、设备慢、设备复杂且共享的问题。DMA/通道减少 CPU 搬运，缓冲平滑速度差，SPOOLing 把独占设备虚拟成共享设备。}
\exam{常考设备分类、四种 I/O 控制方式、DMA 与中断/通道区别、缓冲和 SPOOLing 的作用。}


\concept{循环缓冲与缓冲池}
\begin{points}
  \item 循环缓冲用多个缓冲区构成环形队列，适合连续数据流。
  \item 缓冲池由系统统一管理多个缓冲区，按需分配和回收。
  \item 缓冲池能提高缓冲区利用率，适合多设备、多进程共享。
\end{points}
\plain{I/O 部分就是解决 CPU 快、设备慢、设备复杂且共享的问题。DMA/通道减少 CPU 搬运，缓冲平滑速度差，SPOOLing 把独占设备虚拟成共享设备。}
\exam{常考设备分类、四种 I/O 控制方式、DMA 与中断/通道区别、缓冲和 SPOOLing 的作用。}


\concept{缓冲与缓存区别}
\begin{points}
  \item 缓冲 buffer 主要解决速度不匹配和数据传输粒度不一致。
  \item 缓存 cache 主要利用局部性保存数据副本，减少慢速设备访问。
  \item 例子：磁盘块缓存是 cache；键盘输入缓冲区更偏 buffer。
\end{points}
\plain{I/O 部分就是解决 CPU 快、设备慢、设备复杂且共享的问题。DMA/通道减少 CPU 搬运，缓冲平滑速度差，SPOOLing 把独占设备虚拟成共享设备。}
\exam{常考设备分类、四种 I/O 控制方式、DMA 与中断/通道区别、缓冲和 SPOOLing 的作用。}


\section{设备分配与 SPOOLing}
\concept{设备分配中的数据结构}
\begin{points}
  \item 设备控制表 DCT：描述一台设备的状态、类型、队列等。
  \item 控制器控制表 COCT：描述设备控制器状态。
  \item 通道控制表 CHCT：描述通道状态。
  \item 系统设备表 SDT：记录系统中所有设备及其控制表入口。
\end{points}
\plain{I/O 部分就是解决 CPU 快、设备慢、设备复杂且共享的问题。DMA/通道减少 CPU 搬运，缓冲平滑速度差，SPOOLing 把独占设备虚拟成共享设备。}
\exam{常考设备分类、四种 I/O 控制方式、DMA 与中断/通道区别、缓冲和 SPOOLing 的作用。}


\concept{设备分配过程}
\begin{points}
  \item 根据用户请求的逻辑设备名查找系统设备表。
  \item 分配设备、控制器和通道，若某一级不可用则进程等待。
  \item 分配成功后启动 I/O 操作，I/O 完成后释放相关资源并唤醒等待进程。
\end{points}
\plain{I/O 部分就是解决 CPU 快、设备慢、设备复杂且共享的问题。DMA/通道减少 CPU 搬运，缓冲平滑速度差，SPOOLing 把独占设备虚拟成共享设备。}
\exam{常考设备分类、四种 I/O 控制方式、DMA 与中断/通道区别、缓冲和 SPOOLing 的作用。}


\concept{SPOOLing 系统}
\begin{points}
  \item SPOOLing 是外部设备联机并行操作技术，常用于把独占设备改造成共享设备。
  \item 基本思想：把慢速独占设备的输入/输出先放到磁盘上的输入井/输出井中，由后台进程统一和实际设备交互。
  \item 打印机例子：用户进程把打印数据写入磁盘输出井后即可继续执行，打印进程再按队列送往物理打印机。
  \item SPOOLing 组成：输入井、输出井、输入进程、输出进程、井管理程序。
  \item 直觉：用户不是直接占用打印机，而是把任务交到“打印队列”；物理打印机仍独占，但逻辑上多个用户可同时提交打印任务。
\end{points}
\plain{I/O 部分就是解决 CPU 快、设备慢、设备复杂且共享的问题。DMA/通道减少 CPU 搬运，缓冲平滑速度差，SPOOLing 把独占设备虚拟成共享设备。}
\exam{常考设备分类、四种 I/O 控制方式、DMA 与中断/通道区别、缓冲和 SPOOLing 的作用。}


